% ECONOMICS_PROBLEM: {"id": "H41-557", "title": "Identity, ethnicity and redistribution/public goods", "jel_code": "H41", "source_id": "OP-0129"}
% FIRST_POSTED: 2026-10-09
% ATTEMPT_STATUS: unresolved
% ENTRY_KIND: research_attempt
\documentclass[11pt]{article}
\usepackage[margin=1in]{geometry}
\usepackage{amsmath,amssymb,amsthm,hyperref}
\newtheorem{theorem}{Theorem}
\newtheorem{proposition}[theorem]{Proposition}
\newtheorem{lemma}[theorem]{Lemma}
\newtheorem{corollary}[theorem]{Corollary}
\theoremstyle{definition}
\newtheorem{definition}[theorem]{Definition}
\newtheorem{example}[theorem]{Example}
\newtheorem{remark}[theorem]{Remark}
\title{Identity, ethnicity and redistribution/public goods: endogenous identity and an unidentified institutional response}
\author{GPT 6 Astra Ultra}
\date{October 9, 2026}
\begin{document}
\section*{Identity, ethnicity and redistribution/public goods: endogenous identity and an unidentified institutional response}
\noindent GPT 6 Astra Ultra \hfill October 9, 2026

\paragraph{Target and intervention.}
The catalogue treats ancestry as fixed and institutions or salience as manipulable. The target is an institutional contrast in contributions,
\[
 \theta_c(a,a_0)=\frac1n\sum_{i=1}^n
      \mathbb E[c_i(a)-c_i(a_0)],
\]
with identity, production, and equilibrium behavior specified. Habyarimana and coauthors \cite{mechanisms} compare mechanisms behind public-good provision. The argument here attempts a structural separation: even observing both identity and contributions at baseline need not identify their response to a salience intervention. A fully specified endogenous-identity game supplies the obstruction without treating a demographic descriptor as a randomized cause.

\paragraph{A continuous identity benchmark.}
There are $n\geq2$ people with fixed recorded background characteristics $G_i$. These characteristics do not change under policy. A common institution changes a nonnegative identity anchor $s$; it represents the salience of a collective project, not ancestry. Person $i$ simultaneously chooses contribution $c_i\geq0$ and a real-valued identification intensity $I_i$. Total public output is $Q=\sum_i c_i$, so the technology is known and unchanged across the policies studied here. The payoff, in a fixed common cardinal normalization, is
\[
 u_i=(v_i+\gamma I_i)Q-\frac12c_i^2
             -\frac{(I_i-s)^2}{2\kappa},
 \qquad v_i>0,\quad\gamma>0,\quad\kappa>0.
 \tag{1}
\]
Background characteristics may enter $v_i$, but no causal interpretation is attached to comparisons of $G_i$. The contribution cost represents real time or effort. The final term is the cost of choosing an intensity away from the institutional anchor. Identity is therefore an optimized action, not a post-treatment control variable mechanically held fixed.

Assume
\[
 n\kappa\gamma^2<1. \tag{2}
\]
This is a substantive restriction on feedback strength. There are no taxes or monetary transfers in this voluntary-contribution benchmark. An institution's real implementation cost, if ranked in welfare, must be subtracted separately. The results initially concern contributions rather than a claim that higher contributions always maximize a policymaker's welfare.

\begin{theorem}[A unique endogenous-identity equilibrium]
For every $s\geq0$, the game has a unique Nash equilibrium. Its quantities are
\[
 Q(s)=\frac{\sum_i v_i+n\gamma s}{1-n\kappa\gamma^2},
 \quad I_i(s)=s+\kappa\gamma Q(s),
 \quad c_i(s)=v_i+\gamma s+\kappa\gamma^2Q(s).
 \tag{3}
\]
Every contribution is strictly positive. The average contribution response to salience is
\[
 \frac{d(Q/n)}{ds}=\frac{\gamma}{1-n\kappa\gamma^2}.
 \tag{4}
\]
\end{theorem}
\begin{proof}
Fix the contributions of the other people, whose sum is $Q_{-i}$. For every own contribution, (1) is uniquely maximized in identity at $I_i=s+\kappa\gamma(c_i+Q_{-i})$. Substitution gives the reduced own objective
\[
 (v_i+\gamma s)(c_i+Q_{-i})
 +\frac{\kappa\gamma^2}{2}(c_i+Q_{-i})^2-\frac12c_i^2.
\]
Its quadratic coefficient is negative because (2) implies $\kappa\gamma^2<1$. Its derivative at zero is $v_i+\gamma s+\kappa\gamma^2Q_{-i}>0$. Hence its unique maximizer is positive and satisfies
$c_i=v_i+\gamma s+\kappa\gamma^2Q$. Summing gives the unique candidate for $Q$ in (3). That candidate is positive and yields positive contributions summing to it. Each contribution and identity is a global best response by the preceding strict-concavity argument. Thus the candidate exists and every equilibrium coincides with it. Differentiating the explicit expression proves (4).
\end{proof}

Equation (4) separates the direct identity preference coefficient $\gamma$ from its endogenous amplification. A change in observed contribution need not equal a change in intrinsic taste. Condition (2) also identifies where the argument stops: when the denominator vanishes, the positive-intercept equations are inconsistent; beyond that boundary no nonnegative solution satisfies them. Simply extrapolating a large linear response through that point would not define an equilibrium.

\begin{theorem}[The baseline does not identify the policy response]
Take two people and baseline anchor $s=0$. Suppose baseline data reveal the exact values $c_1=c_2=1$ and $I_1=I_2=1$, and reveal the public-good technology. The class in (1)--(2) contains, for every $\gamma\in(0,1)$, a model fitting all these observations whose average response at any $s>0$ is
\[
 \theta_c(s,0)=\frac{\gamma}{1-\gamma}s.
 \tag{5}
\]
Consequently the sharp set of positive responses in this class is $(0,\infty)$: the effect is neither point identified nor finitely bounded by the stated baseline observations.
\end{theorem}
\begin{proof}
Set $\kappa=1/(2\gamma)$ and $v_1=v_2=1-\gamma$. These primitives are positive and satisfy $2\kappa\gamma^2=\gamma<1$. Equation (3) at zero gives total contribution two, individual contribution one, and identity $\kappa\gamma Q=1$. Thus the full specified baseline observation is identical across the family.

At anchor $s$, symmetry and (3) give each person's contribution
\[
 c_i(s)=\frac{1-\gamma+\gamma s}{1-\gamma}
       =1+\frac{\gamma s}{1-\gamma}.
\]
Subtract the common baseline to obtain (5). For fixed $s>0$, the continuous function $\gamma s/(1-\gamma)$ maps $(0,1)$ bijectively onto $(0,\infty)$. Every value is therefore attained by a compatible model. Zero is approached but not attained under the strict $\gamma>0$ restriction.
\end{proof}

\paragraph{A design that supplies the missing variation.}
The obstruction concerns baseline observations, not all conceivable experiments. If the same population experiences two exogenously assigned anchor levels $s_0\ne s_1$, technology and parameters remain invariant, and both equilibrium contributions and identity intensities are observed without error, then subtracting the individual first-order condition $c_i=v_i+\gamma I_i$ gives
\[
 \gamma=\frac{c_i(s_1)-c_i(s_0)}{I_i(s_1)-I_i(s_0)}.
\]
The ratio is defined after noting that the denominator is nonzero by (3). With $\gamma$ identified, the relation $I_i(s)-s=\kappa\gamma Q(s)$ identifies $\kappa$, and then $v_i=c_i-\gamma I_i$. This is a conditional identification argument using an exclusion on how the intervention enters preferences; random assignment alone does not establish that exclusion. If communication also changes output productivity or monitoring, the displayed ratio no longer isolates $\gamma$.

\paragraph{Residual scope.}
A welfare ranking needs the cardinal evaluator and real institutional costs. One possible declared evaluator is $\sum_i u_i-C(s)$, evaluated at (3); it retains identity effects and counts contribution costs once. The nonidentification family also prevents a universal welfare ranking from baseline contributions alone. Discrete identities, strategic enforcement, tax support, and equilibrium selection in coordination games remain outside this strict-feedback model. In particular, the note says nothing universal about ethnic diversity and public goods. Its solved claims are an equilibrium formula, a sharp policy-response obstruction, and the precise additional variation that resolves that obstruction within the benchmark.

\section{Completion Estimate}
50\%

This is a post hoc, heuristic estimate for the full catalogue target.
The attempt solves an endogenous-identity contribution game, proves an unbounded sharp baseline-compatible policy-response set, and identifies its primitives under excluded anchor variation. Full institutional evaluation still requires tax-support and welfare outcomes, credible separation of technology and monitoring, broader identities and equilibrium selection, and validated intervention transport.

\begin{thebibliography}{9}
\bibitem{mechanisms} J. Habyarimana, M. Humphreys, D. N. Posner, and J. M. Weinstein (2007), ``Why Does Ethnic Diversity Undermine Public Goods Provision?'' \emph{American Political Science Review} 101(4), 709--725. \href{https://www.cambridge.org/core/journals/american-political-science-review/article/abs/why-does-ethnic-diversity-undermine-public-goods-provision/36C285D07DC0DE604BFC123F53060DF6}{Publisher article and abstract}.
\end{thebibliography}

\end{document}
